\providecommand{\Ric}{\operatorname{Ric}}
\providecommand{\divergence}{\operatorname{div}}
\providecommand{\grad}{\operatorname{grad}}
\providecommand{\Exp}{\operatorname{exp}}
\providecommand{\Log}{\operatorname{log}}
\providecommand{\TpM}{T_p\mathcal{M}}
\providecommand{\M}{\mathcal{M}}

\chapter{Geodesics: Learning on Manifolds}
\label{ch:geodesics}

The previous chapters built two pillars of geometric deep learning: the
\emph{domains} where data live and the \emph{graphs} that discretize them.
\Cref{ch:graphs} showed, in particular, that the Laplacian of a proximity graph
converges to the Laplace--Beltrami operator of the underlying surface, and used
this connection to define intrinsic Fourier transforms and spectral filters.
That operator, however, presupposes that the domain carries a \emph{Riemannian
metric} --- a notion the smooth manifolds of \Cref{ch:foundations} do not yet
possess: they have topological and differentiable structure, but no intrinsic
sense of distance.

\emph{Geodesics} fill this gap. In Euclidean space the shortest path between two
points is trivially a straight line; on a curved manifold, geodesics generalize
this idea as curves that locally minimize length. They are the thread that ties
the preceding themes together: they supply the intrinsic distances that justify
the weights of proximity graphs, they define the local coordinate systems needed
for geometric convolutions, and they govern how information propagates in a way
compatible with the curvature of the domain. To make these ideas precise we
first endow smooth manifolds with metric structure
(\Cref{sec:geo:intrinsic}), introduce the connection and curvature that measure
the departure from flatness, define geodesics together with the exponential and
logarithmic maps (\Cref{sec:geo:geodesics}), and then turn to learning:
geodesic convolution (\Cref{sec:geo:convolution}) and spectral methods on
manifolds (\Cref{sec:geo:spectral}).

\section{Intrinsic geometry}
\label{sec:geo:intrinsic}

\subsection{Riemannian metrics}
\label{subsec:geo:metrics}

To speak of geodesics rigorously we must equip the differentiable manifolds of
\Cref{ch:foundations} with metric structure. A \emph{Riemannian metric}
provides the machinery to measure lengths, angles, and volumes intrinsically ---
that is, using only quantities available on the manifold itself, without
reference to an ambient space.

\begin{definition}[Riemannian metric]
\label{def:geo:metric}
A \emph{Riemannian metric} on a differentiable manifold $\M$ is a smooth
assignment to each point $p \in \M$ of an inner product
$g_p \colon \TpM \times \TpM \to \R$ on the tangent space. In local coordinates
$(x^1, \dots, x^m)$ it is written
\[
  g_p = \sum_{i,j=1}^m g_{ij}(p)\, dx^i \otimes dx^j ,
\]
where the $g_{ij}(p)$ are smooth functions forming a symmetric positive-definite
matrix at every point. The pair $(\M, g)$ is a \emph{Riemannian manifold}.
\end{definition}

\begin{example}[Metrics on common manifolds]
\label{ex:geo:metrics}
\leavevmode
\begin{itemize}
  \item \textbf{Euclidean space} $\R^n$: the standard metric is $g_{ij} =
    \delta_{ij}$, giving $ds^2 = dx_1^2 + \cdots + dx_n^2$.
  \item \textbf{Sphere} $S^2$: in spherical coordinates $(\theta, \phi)$, the
    metric induced by the embedding in $\R^3$ is
    $ds^2 = d\theta^2 + \sin^2\theta\, d\phi^2$.
  \item \textbf{Hyperbolic space} $\mathbb{H}^n$ (Poincar\'e model): on the disc
    $\mathbb{D}^n = \{x \in \R^n : \norm{x} < 1\}$ the metric is
    $g_{ij} = 4\delta_{ij} / (1 - \norm{x}^2)^2$, which ``stretches'' distances
    near the boundary.
\end{itemize}
\end{example}

The metric supplies everything needed to measure length and distance, through a
three-step chain. First, since $g_p$ is an inner product on each tangent space,
it induces a pointwise norm $\norm{v}_g = \sqrt{g_p(v,v)}$; for a curve
$\gamma(t)$, this measures the instantaneous speed $\norm{\dot\gamma(t)}_g$.
Second, integrating the speed along the curve yields its total length. Third,
taking the infimum of length over all curves joining two points gives the
geodesic distance.

\begin{definition}[Length and geodesic distance]
\label{def:geo:distance}
The \emph{length} of a curve $\gamma \colon [a,b] \to \M$ is
\[
  L(\gamma) = \int_a^b \sqrt{g_{\gamma(t)}\big(\dot\gamma(t), \dot\gamma(t)\big)}\, dt ,
\]
and the \emph{geodesic distance} between $p, q \in \M$ is
\[
  d_g(p,q) = \inf\{ L(\gamma) : \gamma \text{ a smooth curve from } p \text{ to } q\} .
\]
\end{definition}

\begin{proposition}[The geodesic distance is a metric]
\label{prop:geo:dg-metric}
On a connected Riemannian manifold $(\M, g)$, the function $d_g$ is a metric in
the topological sense: for all $p, q, r \in \M$,
$d_g(p,q) \ge 0$ with equality iff $p = q$; $d_g(p,q) = d_g(q,p)$; and
$d_g(p,r) \le d_g(p,q) + d_g(q,r)$.
\end{proposition}

Symmetry follows from reversing curve parametrizations, and the triangle
inequality from concatenating near-optimal curves. Non-degeneracy is the only
delicate point: if $p \neq q$, a compactness argument on a coordinate ball shows
$g_x(v,v) \ge c\,\norm{v}_{\mathrm{euc}}^2$ for some $c > 0$, so any curve
leaving the ball has length bounded below by a positive constant independent of
the curve.

\begin{example}[Geodesic distance on the sphere]
\label{ex:geo:sphere-distance}
On the unit sphere $S^2$ with metric $ds^2 = d\theta^2 + \sin^2\theta\, d\phi^2$,
the great-circle arc from the north pole $p$ to
$q = (\sin\alpha, 0, \cos\alpha)$ is $\gamma(t) = (\sin t, 0, \cos t)$,
$t \in [0,\alpha]$. Its speed is $g_{\gamma(t)}(\dot\gamma, \dot\gamma) = 1$, so
$L(\gamma) = \alpha$. Since great-circle arcs minimize length,
$d_g(p,q) = \alpha = \arccos\inner{p}{q}$, with $\inner{\cdot}{\cdot}$ the
Euclidean inner product of $\R^3$.
\end{example}

The metric also yields angles, $\cos\theta = g_p(v,w) /
(\norm{v}_g \norm{w}_g)$, and an intrinsic \emph{volume element}
$dV_g = \sqrt{\det(g_{ij})}\, dx^1 \wedge \cdots \wedge dx^m$, which lets one
integrate functions over the manifold.

\subsection{The Laplace--Beltrami operator}
\label{subsec:geo:laplacian}

The Euclidean Laplacian $\nabla^2 f = \sum_i \partial^2 f / \partial(x^i)^2 =
\divergence(\grad f)$ measures how the value of $f$ at a point deviates from its
local average. Replacing the Euclidean gradient and divergence by their
metric counterparts yields its manifold generalization.

\begin{definition}[Laplace--Beltrami operator]
\label{def:geo:laplace-beltrami}
On a Riemannian manifold $(\M, g)$, the \emph{Laplace--Beltrami operator}
$\Delta_\M$ acts on smooth functions $f \colon \M \to \R$ as
\[
  \Delta_\M f = \divergence(\grad f)
    = \frac{1}{\sqrt{\det g}}\, \partial_i\!\left(\sqrt{\det g}\; g^{ij}\, \partial_j f\right),
\]
where $g^{ij}$ are the components of the inverse metric. When
$(\M, g) = (\R^m, \delta_{ij})$ one has $\det g = 1$ and $g^{ij} = \delta^{ij}$,
so $\Delta_\M$ reduces to the Euclidean Laplacian.
\end{definition}

The Laplace--Beltrami operator is the single object that links differential
geometry to geometric deep learning, for four reasons.

\begin{remark}[Why the Laplace--Beltrami operator matters]
\label{rem:geo:lb-role}
\leavevmode
\begin{enumerate}
  \item \textbf{Diffusion.} $\Delta_\M$ governs the heat equation
    $\partial_t u = \Delta_\M u$, which models the spread of information across
    the manifold. A layer that ``diffuses'' features toward geodesic
    neighbourhoods is the learnable analogue of this process.
  \item \textbf{Intrinsic spectral analysis.} The eigenfunctions
    $\{\varphi_k\}$ of $\Delta_\M$, with $\Delta_\M \varphi_k = \lambda_k
    \varphi_k$, form an orthonormal basis of $L^2(\M)$ that generalizes the
    Fourier basis, enabling an intrinsic Fourier transform
    $\hat f(\lambda_k) = \inner{f}{\varphi_k}$ and hence spectral convolutions
    (\Cref{sec:geo:spectral}).
  \item \textbf{Isometry invariance.} For every isometry $\phi$ of $\M$,
    $\Delta_\M(f \circ \phi) = (\Delta_\M f) \circ \phi$. Any architecture built
    from $\Delta_\M$ therefore inherits invariance under the manifold's
    symmetries --- exactly the equivariance principle of \Cref{ch:priors}.
  \item \textbf{Discretization.} The graph Laplacian of \Cref{ch:graphs} is the
    discrete analogue of $\Delta_\M$ and, under suitable sampling, converges to
    it, so techniques transfer between the continuous and discrete worlds.
\end{enumerate}
\end{remark}

\subsection{Connections and curvature}
\label{subsec:geo:curvature}

To define geodesics and quantify how a manifold departs from flat space we need
two more tools. In $\R^n$ all tangent spaces are canonically identified with
$\R^n$, so vectors at different points can be compared directly; on a general
manifold $\TpM$ and $T_q\M$ are distinct vector spaces with no canonical
identification. A \emph{connection} repairs this by prescribing how to
differentiate vector fields, and the \emph{curvature tensor} then measures the
non-commutativity of the resulting covariant derivatives.

\begin{definition}[Affine connection]
\label{def:geo:connection}
An \emph{affine connection} on $\M$ is a map $\nabla \colon \Gamma(T\M) \times
\Gamma(T\M) \to \Gamma(T\M)$, $(X,Y) \mapsto \nabla_X Y$ on smooth vector fields
that is $C^\infty(\M)$-linear in $X$, $\R$-linear in $Y$, and satisfies the
Leibniz rule $\nabla_X(fY) = (Xf)Y + f\nabla_X Y$.
\end{definition}

A manifold admits infinitely many connections, but a metric singles out a
canonical one.

\begin{theorem}[Fundamental theorem of Riemannian geometry]
\label{thm:geo:levi-civita}
On a Riemannian manifold $(\M, g)$ there is a unique affine connection $\nabla$,
the \emph{Levi--Civita connection}, that is
\emph{metric-compatible}, $X g(Y,Z) = g(\nabla_X Y, Z) + g(Y, \nabla_X Z)$, and
\emph{torsion-free}, $\nabla_X Y - \nabla_Y X = [X,Y]$.
In local coordinates it is determined by the \emph{Christoffel symbols}
\[
  \Gamma^k_{ij} = \tfrac{1}{2} g^{kl}\!\left(
    \partial_j g_{il} + \partial_i g_{jl} - \partial_l g_{ij}\right).
\]
\end{theorem}

Uniqueness follows from the \emph{Koszul formula}, obtained by cyclically
expanding metric compatibility and using torsion-freeness to cancel the cross
terms; since $g$ is non-degenerate, the formula determines $\nabla_X Y$
uniquely, and a direct check shows it satisfies the axioms. Geometrically, the
two defining properties say that parallel transport preserves lengths and
angles, and that infinitesimal parallelograms close --- transport introduces no
spurious rotation.

The Levi--Civita connection lets us compare the second covariant derivatives
$\nabla_X \nabla_Y Z$ and $\nabla_Y \nabla_X Z$. In Euclidean space they agree;
on a curved manifold they differ, and the discrepancy is the curvature.

\begin{definition}[Riemann curvature tensor]
\label{def:geo:riemann}
The \emph{Riemann curvature tensor} of $(\M, g)$ is
\[
  R(X,Y)Z = \nabla_X \nabla_Y Z - \nabla_Y \nabla_X Z - \nabla_{[X,Y]} Z .
\]
It is $C^\infty(\M)$-multilinear, so its value at $p$ depends only on $X, Y, Z$
at $p$. A Riemannian manifold is locally isometric to $\R^n$ if and only if
$R \equiv 0$.
\end{definition}

\begin{definition}[Sectional, Ricci, and scalar curvature]
\label{def:geo:sectional}
For a $2$-plane $\sigma \subset \TpM$ spanned by independent $u, v$, the
\emph{sectional curvature} is
\[
  K(\sigma) = \frac{\inner{R(u,v)v}{u}_g}{\norm{u}^2\norm{v}^2 - \inner{u}{v}_g^2},
\]
the curvature of the surface swept out by geodesics in directions of $\sigma$.
Contracting $R$ gives the \emph{Ricci tensor}
$R_{ij} = \sum_k \inner{R(\partial_i,\partial_k)\partial_j}{\partial_k}_g$, and
tracing again gives the \emph{scalar curvature} $\mathcal{R} = g^{ij}R_{ij}$.
\end{definition}

Sectional curvature controls how geometry departs from Euclid's postulates
(\Cref{fig:geo:curvature}). In a geodesic triangle of area $A$ the angle sum is
$\alpha + \beta + \gamma = \pi + K\,A$; a geodesic disc of radius $r$ has area
growing faster than $\pi r^2$ when $K < 0$ and slower when $K > 0$; and the
separation $\xi(t)$ between initially parallel geodesics obeys the Jacobi
equation $\ddot\xi + K\xi = 0$, whose solutions are oscillatory ($K > 0$),
constant ($K = 0$), or exponential ($K < 0$). This last fact explains why
hyperbolic space, with its exponential volume growth, embeds trees with low
distortion --- a property we exploit in \Cref{subsec:geo:specific}.

\begin{figure}[h]
\centering
\begin{tikzpicture}[scale=1.0]
% --- Negative curvature (saddle) ---
\begin{scope}[shift={(-5.2,0)}]
  \shade[ball color=gdlred!18, opacity=0.85]
    (-1.6,0.6) .. controls (-0.6,1.5) and (0.6,-0.5) .. (1.6,0.6)
    .. controls (1.0,-0.2) and (0.4,-1.2) .. (0,-1.0)
    .. controls (-0.4,-1.2) and (-1.0,-0.2) .. (-1.6,0.6);
  \draw[gdlred!50, thin] (-1.4,0.5) .. controls (0,1.0) .. (1.4,0.5);
  \draw[gdlred!50, thin] (-0.9,-0.7) .. controls (0,-0.3) .. (0.9,-0.7);
  \node[below, font=\small\bfseries, text=gdlred!70!black] at (0,-1.6) {$K<0$};
  \node[below, font=\footnotesize, text=gdlred!70!black] at (0,-2.0) {diverge};
\end{scope}
% --- Zero curvature (plane) ---
\begin{scope}[shift={(0,0)}]
  \fill[gdlgreen!15] (-1.6,-1.0) -- (1.6,-1.0) -- (1.6,1.0) -- (-1.6,1.0) -- cycle;
  \foreach \x in {-1.2,-0.6,0,0.6,1.2}
    \draw[gdlgreen!45, thin] (\x,-1.0) -- (\x,1.0);
  \foreach \y in {-0.6,0,0.6}
    \draw[gdlgreen!45, thin] (-1.6,\y) -- (1.6,\y);
  \node[below, font=\small\bfseries, text=gdlgreen!60!black] at (0,-1.6) {$K=0$};
  \node[below, font=\footnotesize, text=gdlgreen!60!black] at (0,-2.0) {parallel};
\end{scope}
% --- Positive curvature (sphere) ---
\begin{scope}[shift={(5.2,0)}]
  \shade[ball color=gdlblue!22, opacity=0.9] (0,0) circle (1.3cm);
  \draw[gdlblue!50, thin] (0,0) circle (1.3cm);
  \draw[gdlblue!50, thin] (-1.3,0) .. controls (0,0.35) .. (1.3,0);
  \draw[gdlblue!50, thin] (0,1.3) .. controls (0.35,0) .. (0,-1.3);
  \node[below, font=\small\bfseries, text=gdlblue!70!black] at (0,-1.6) {$K>0$};
  \node[below, font=\footnotesize, text=gdlblue!70!black] at (0,-2.0) {converge};
\end{scope}
\end{tikzpicture}
\caption[Curvature comparison]{The three regimes of sectional curvature.
Geodesics that start parallel diverge on a negatively curved surface
($K<0$, left), stay parallel on a flat plane ($K=0$, centre), and converge on a
positively curved one ($K>0$, right).}
\label{fig:geo:curvature}
\end{figure}

\section{Geodesics, exponential, and logarithmic maps}
\label{sec:geo:geodesics}

\emph{Geodesics} are the curves that locally minimize length, generalizing
straight lines to curved space. The Levi--Civita connection supplies the
covariant derivative that formalizes them: a geodesic is a curve whose tangent
vector is parallel along itself, so it ``does not accelerate'' relative to the
geometry.

\begin{definition}[Geodesic]
\label{def:geo:geodesic}
A curve $\gamma \colon I \to \M$ is a \emph{geodesic} if
$\nabla_{\dot\gamma}\dot\gamma = 0$, which in local coordinates is the
second-order system
\[
  \ddot\gamma^k + \Gamma^k_{ij}\, \dot\gamma^i \dot\gamma^j = 0 ,
\]
using the Einstein summation convention.
\end{definition}

Geodesics admit three complementary readings: as extremals of the length
functional (variational principle), as the trajectories of a free particle
(inertia), and as locally shortest paths. They are straight lines in $\R^n$,
great circles on $S^2$, and arcs of circles meeting the boundary orthogonally in
the Poincar\'e model of $\mathbb{H}^n$.

\begin{definition}[Exponential and logarithmic maps]
\label{def:geo:exp-log}
The \emph{exponential map} $\exp_p \colon \TpM \to \M$ sends $v$ to
$\gamma_v(1)$, where $\gamma_v$ is the unique geodesic with $\gamma_v(0) = p$ and
$\dot\gamma_v(0) = v$ --- the point reached by walking the geodesic in direction
$v$ for unit time. Its local inverse, the \emph{logarithmic map}
$\log_p \colon U \to \TpM$, is defined on a neighbourhood $U$ of $p$ where
$\exp_p$ is a diffeomorphism: $\log_p(q)$ is the ``vector pointing from $p$ to
$q$.''
\end{definition}

\begin{proposition}[Properties of the exponential map]
\label{prop:geo:exp}
For $p \in \M$: (i) $\exp_p(0) = p$; (ii) $\exp_p(tv) = \gamma_v(t)$; (iii)
$d(\exp_p)_0 = \operatorname{id}_{\TpM}$; (iv) $\exp_p$ is a diffeomorphism from
a neighbourhood of $0 \in \TpM$ onto a neighbourhood of $p$; and (v)
$d_g(p, \exp_p(v)) = \norm{v}_g$ for $v$ small enough.
\end{proposition}

Property (iii) follows from differentiating $t \mapsto \exp_p(tv)$ at $t = 0$;
the inverse function theorem then gives (iv). Property (v) holds because the
Levi--Civita connection keeps $\norm{\dot\gamma_v}_g$ constant, so the geodesic
$\gamma_v|_{[0,1]}$ has length $\norm{v}_g$ and is minimizing for small $v$. The
identity $d_g(p,q) = \norm{\log_p(q)}_g$ is exactly the bridge between intrinsic
distance and the tangent space (\Cref{fig:geo:exp-log}).

\begin{figure}[h]
\centering
\begin{tikzpicture}[>=Stealth, scale=1.0]
% Manifold (sphere)
\shade[ball color=gdlblue!18, opacity=0.9] (0,0) circle (2.6cm);
\draw[thick, gdlblue!45] (0,0) circle (2.6cm);
\node[gdlblue!60!black, font=\small] at (-2.1,-2.0) {$\mathcal{M}$};
% Point p
\coordinate (p) at (0,2.6);
\fill[gdlred!75] (p) circle (2.5pt);
\node[above right, font=\small] at (p) {$p$};
% Tangent plane
\begin{scope}[shift={(0,2.6)}]
  \fill[gdlgreen!16] (-2.2,-0.25) -- (2.2,-0.25) -- (2.2,1.4) -- (-2.2,1.4) -- cycle;
  \draw[gdlgreen!50] (-2.2,0) -- (2.2,0);
  \node[gdlgreen!55!black, font=\small] at (2.4,1.0) {$T_p\mathcal{M}$};
  \draw[->, very thick, gdlorange!80] (0,0) -- (1.3,0.7) node[above, font=\small] {$v$};
\end{scope}
% Geodesic
\draw[very thick, gdlred!70, dashed] (0,2.6) arc (90:33:2.6cm);
% Point q
\coordinate (q) at ({2.6*cos(33)},{2.6*sin(33)});
\fill[gdlpurple!75] (q) circle (2.5pt);
\node[right, font=\small, gdlpurple!75] at (q) {$q=\exp_p(v)$};
% Maps
\draw[->, very thick, gdlorange!80] (1.3,3.3) .. controls (2.3,2.8) and (2.6,2.0) .. (q);
\node[gdlorange!80, font=\small] at (3.0,3.0) {$\exp_p$};
\draw[->, very thick, gdlpurple!75, dashed] ($(q)+(-0.25,0.3)$) .. controls (2.2,2.5) and (1.7,3.4) .. (1.0,3.5);
\node[gdlpurple!75, font=\small] at (2.6,3.7) {$\log_p$};
% Note
\node[font=\footnotesize, text width=4.4cm, align=left] at (-4.6,0.2) {
  $\exp_p(v)=\gamma_v(1)$, $\gamma_v$ geodesic with $\dot\gamma_v(0)=v$\\[2pt]
  $\log_p=\exp_p^{-1}$ (local)\\[2pt]
  $d_g(p,q)=\norm{\log_p(q)}_g$};
\end{tikzpicture}
\caption[Exponential and logarithmic maps]{The exponential map sends a tangent
vector $v \in T_p\mathcal{M}$ to the endpoint of the geodesic it generates; the
logarithmic map is its local inverse. Geodesic distance equals the tangent-space
norm of $\log_p(q)$.}
\label{fig:geo:exp-log}
\end{figure}

\paragraph{Normal coordinates.}
Choosing an orthonormal basis $\{e_1, \dots, e_m\}$ of $\TpM$ and setting
$\phi(\exp_p(v^i e_i)) = (v^1, \dots, v^m)$ defines \emph{geodesic normal
coordinates} centred at $p$. In these coordinates geodesics through $p$ are
straight lines through the origin, $g_{ij}(p) = \delta_{ij}$, the Christoffel
symbols vanish at $p$, and so do the first derivatives of the metric. The
second-order term reveals the curvature:
\[
  g_{ij}(x) = \delta_{ij} - \tfrac{1}{3} R_{ikjl}(p)\, x^k x^l + O(\norm{x}^3) .
\]
Normal coordinates thus \emph{linearize} the manifold: within a geodesic ball of
radius $r$ it looks Euclidean up to $O(r^2)$ corrections governed by curvature.
This is precisely the chart in which geodesic convolutions operate.

\begin{remark}[The gauge ambiguity]
\label{rem:geo:gauge}
The basis $\{e_i\}$ is not canonical: any other orthonormal basis differs by a
rotation $O \in O(m)$. This \emph{gauge freedom} rotates the normal-coordinate
chart, so a convolution kernel expressed in these coordinates transforms as
$k'(r,\theta) = k(r, \theta - \alpha)$. Making a network independent of the
choice requires equivariance constraints on the kernel --- the subject of
\Cref{ch:gauges}.
\end{remark}

\section{Geodesic convolution}
\label{sec:geo:convolution}

Group convolutions (\Cref{ch:groups}) exploit a global symmetry, but many
real-world domains --- 3D surfaces, cortical sheets, configuration spaces ---
carry no such symmetry: their local geometry varies from point to point, there
is no globally consistent translation, and no canonical global coordinate
system. The exponential map removes these obstacles locally, and \emph{geodesic
convolutional neural networks}~\citep{masci2015geodesic} were the first
systematic extension of convolution to general Riemannian manifolds.

\subsection{Convolution in geodesic charts}
\label{subsec:geo:gcnn}

The idea is to build, around each point $p$, a local chart via the exponential
map and to perform a polar-coordinate convolution inside it. A signal
$f \colon \M \to \R$ is convolved with a learnable kernel $k$ by sampling the
manifold along geodesic rays:
\begin{equation}
\label{eq:geo:gcnn}
  (f *_g k)(p) = \sum_{i,j} k(r_i, \theta_j)\;
    f\!\left(\exp_p(r_i\, e_{\theta_j})\right)\; J_p(r_i, \theta_j),
\end{equation}
where $e_{\theta_j} = \cos\theta_j\, e_1 + \sin\theta_j\, e_2$ is the angular
direction, $(r_i, \theta_j)$ discretizes a geodesic disc in polar coordinates,
and $J_p(r, \theta) = \sqrt{\det g_{\exp_p(r e_\theta)}}$ is a Jacobian factor
correcting the distortion between the tangent patch and the manifold. By the
normal-coordinate expansion of \Cref{rem:geo:gauge}'s neighbourhood,
\[
  J_p(r, \theta) = 1 - \tfrac{1}{6} K(p)\, r^2 + O(r^3),
\]
so the correction is negligible for small patches and grows with the Gaussian
curvature $K(p)$. Operationally, building each chart requires choosing an
orthonormal tangent basis, mapping a polar grid through $\exp_p$, and
interpolating the feature values at the resulting points
(\Cref{fig:geo:convolution}).

\begin{figure}[h]
\centering
\begin{tikzpicture}[>=Stealth, scale=1.0]
% Curved surface patch
\shade[ball color=gdlorange!20, opacity=0.9]
  (-2.6,0) .. controls (-2.0,1.4) and (2.0,1.4) .. (2.6,0)
  .. controls (2.0,-1.2) and (-2.0,-1.2) .. (-2.6,0);
\draw[gdlorange!55, thick]
  (-2.6,0) .. controls (-2.0,1.4) and (2.0,1.4) .. (2.6,0)
  .. controls (2.0,-1.2) and (-2.0,-1.2) .. (-2.6,0);
\node[gdlorange!60!black, font=\small] at (-2.4,-1.1) {$\mathcal{M}$};
% Center point p
\coordinate (p) at (0,0.25);
% Geodesic polar grid (rings + rays)
\foreach \rr in {0.45,0.9,1.35}
  \draw[gdlblue!55, thin] (p) ellipse ({\rr} and {0.55*\rr});
\foreach \aa in {0,45,90,135,180,225,270,315}
  \draw[gdlblue!40, thin] (p) -- ($(p)+({1.35*cos(\aa)} and {0.55*1.35*sin(\aa)})$);
\fill[gdlred!80] (p) circle (2.5pt);
\node[above, font=\small] at (p) {$p$};
% A neighbour
\coordinate (q) at ($(p)+({1.35*cos(35)} and {0.55*1.35*sin(35)})$);
\fill[gdlgreen!70] (q) circle (2pt);
\node[right, font=\small, gdlgreen!60!black] at (q) {$q_j=\exp_p(r_i e_{\theta_j})$};
% label
\node[font=\footnotesize, text=gdlblue!60!black] at (0,-1.5) {geodesic ball $B_r(p)$};
\end{tikzpicture}
\caption[Geodesic convolution]{Geodesic convolution. A polar grid of radii
$r_i$ and angles $\theta_j$ is laid down in the tangent space at $p$ and mapped
onto the geodesic ball $B_r(p)$ through $\exp_p$; the kernel weights
$k(r_i,\theta_j)$ are applied to the interpolated features, corrected by the
Jacobian $J_p$.}
\label{fig:geo:convolution}
\end{figure}

Geodesic convolutions were a conceptual breakthrough, but they have clear
limitations. The construction depends on a choice of tangent basis at each
point, which may be inconsistent between neighbours (the gauge ambiguity of
\Cref{rem:geo:gauge}); computing geodesic charts and Jacobians is expensive on
high-resolution meshes; and the exponential map requires enough smoothness, so
singular domains are problematic. \emph{Anisotropic diffusion
descriptors}~\citep{Boscaini2016AnisotropicDD} address the related shortcoming
of isotropic spectral descriptors by replacing the heat kernel with a
direction-dependent diffusion $\partial_t u = -\divergence(C \nabla u)$ whose
conductivity tensor $C$ adapts to local geometry, capturing the orientation of
elongated features.

\subsection{A general framework: MoNet}
\label{subsec:geo:monet}

The dependence on an explicit coordinate system motivates a more abstract
formulation. MoNet~\citep{DBLP:journals/corr/MontiBMRSB16} replaces geodesic
charts with \emph{pseudo-coordinates}: a map $u \colon \M \times \M \to \R^d$
that encodes the local geometric relation between a point $x$ and its neighbours,
is local (vanishing outside a neighbourhood), isometry-invariant, and smooth.
The convolution becomes a weighted aggregation over a neighbourhood
$\mathcal{N}(x)$,
\begin{equation}
\label{eq:geo:monet}
  (f * w)(x) = \int_{\mathcal{N}(x)} w_\theta\big(u(x,y)\big)\, f(y)\, dV_g(y),
\end{equation}
where the kernel $w_\theta$ is a learnable mixture of parametric basis functions,
typically Gaussians,
\[
  w_\theta(u) = \sum_{j=1}^J \alpha_j\,
    \exp\!\left(-\tfrac{1}{2}(u - \mu_j)^\top \Sigma_j^{-1} (u - \mu_j)\right),
\]
with learnable weights $\alpha_j$, means $\mu_j$, and covariances $\Sigma_j$. On
a discrete mesh the integral becomes
$(f * w)(x_i) = \sum_{j \in \mathcal{N}(i)} w_\theta(u(x_i,x_j))\, f(x_j)\, A_j$
with local area weights $A_j$.

\begin{theorem}[Universal approximation of MoNet kernels]
\label{thm:geo:monet}
For any continuous compactly supported $k \colon \R^d \to \R$ and any
$\varepsilon > 0$ there exist $J$ and parameters
$\{\alpha_j, \mu_j, \Sigma_j\}_{j=1}^J$ with
$\sup_{u}\, \abs{k(u) - w_\theta(u)} < \varepsilon$.
\end{theorem}

This follows from the density of Gaussian mixtures together with compactness of
the support, and it makes MoNet a genuine generalization: different choices of
$u$ recover earlier models. Taking $u(x,y) = (d_g(x,y), \angle(x,y))$ (geodesic
distance and polar angle) reproduces geodesic
convolutions~\eqref{eq:geo:gcnn}; taking degree-based coordinates
$u(x,y) = (1/\sqrt{\deg x},\, 1/\sqrt{\deg y})$ recovers graph convolutional
networks (\Cref{ch:graphs}); and B-spline coordinates recover spline-based
convolutions. MoNet thus unifies grid, graph, and manifold convolutions under a
single learnable-kernel recipe, and is itself an instance of the message-passing
template of \Cref{ch:graphs}.

\subsection{Parallel transport and specific manifolds}
\label{subsec:geo:specific}

Comparing feature vectors that live in different tangent spaces requires moving
one to the other. The Levi--Civita connection provides \emph{parallel transport}
$\Pi_{p \to q}^\gamma \colon \TpM \to T_q\M$ along the geodesic $\gamma$ joining
$p$ and $q$; it is a linear isometry, preserving norms and angles, and composes
along concatenated paths. The holonomy around closed loops is controlled by
curvature (the Ambrose--Singer theorem identifies the holonomy Lie algebra with
the curvature operators), so parallel transport encodes the full intrinsic
geometry. \emph{Parallel transport convolutions}~\citep{schonsheck2018parallel}
exploit this to define a coordinate-free convolution that transports the kernel
consistently across the surface, sidestepping the gauge ambiguity of geodesic
charts. These ideas connect directly to the gauge-equivariant networks of
\Cref{ch:gauges}.

The exponential and logarithmic maps also let classical building blocks be
restated intrinsically. A \emph{Riemannian residual connection} replaces the
Euclidean update $x_{\ell+1} = x_\ell + F_\ell(x_\ell)$ by
\[
  x_{\ell+1} = \exp_{x_\ell}\!\big(F_\ell(x_\ell)\big),
    \qquad F_\ell(x_\ell) \in T_{x_\ell}\M,
\]
which keeps the state on the manifold, and geodesic interpolation
$\gamma(t) = \exp_p(t \log_p(q))$ generalizes linear interpolation.

When the data have a known geometry one can specialize the manifold. Hierarchical
data --- trees, taxonomies, ontologies --- embed naturally in \emph{hyperbolic
space}: because the volume of a hyperbolic ball grows exponentially with its
radius (the $K < 0$ regime of \Cref{fig:geo:curvature}), a tree whose node count
grows exponentially with depth fits with bounded distortion, impossible in any
Euclidean space of fixed dimension. In the Poincar\'e ball one replaces vector
addition by the M\"obius operation and Euclidean distance by the hyperbolic
distance; \emph{hyperbolic attention
networks}~\citep{gulcehre2019hyperbolic} weight attention by
$\exp(-d_{\mathbb{H}}(x,y)/\tau)$ with the hyperbolic distance $d_{\mathbb{H}}$.
The unifying principle is the same throughout: respect the geometry of the data
space when designing every component, so that equivariance under the isometry
group makes the learned representations intrinsic.

\subsection{Case studies}
\label{subsec:geo:cases}

Three applications show geodesic methods outperforming Euclidean baselines.

\paragraph{3D shape analysis (FAUST).}
The FAUST benchmark contains triangulated meshes of human bodies (about $6890$
vertices each) in varied poses, with strong variations in curvature. Because
geodesic distances are invariant under the (nearly isometric) deformations of a
moving body, geodesic convolutions and MoNet recognize correspondences and poses
far more accurately than a Euclidean network applied to a flat projection,
which is sensitive to the embedding. High-curvature regions such as joints
benefit most, and the Jacobian correction is critical near them.

\paragraph{Omnidirectional ($360^\circ$) images.}
A spherical image projected to a flat (equirectangular) plane suffers severe
polar distortion --- areas near the poles are stretched by an order of
magnitude, breaking the translation invariance a standard convolution assumes.
Treating the image as a signal on $S^2$ and convolving geodesically (for
example, over an icosahedral discretization) keeps distortion almost uniform and
preserves local shapes, improving scene classification and semantic
segmentation.

\paragraph{Medical surfaces (cerebral cortex).}
The cortical surface is a genus-0 manifold with deep folds and extreme curvature
variation, and exhibits large anatomical variability across subjects.
Convolving intrinsically on the cortical surface, rather than on a volumetric or
projected representation, respects the folding pattern and yields anatomically
consistent features for tasks such as parcellation across large cohorts.

\section{Spectral methods on manifolds}
\label{sec:geo:spectral}

Geodesic convolution is an \emph{intrinsic but spatial} construction: it works
in geodesic charts. The spectral viewpoint offers a complementary, fully
coordinate-free alternative, built directly on the Laplace--Beltrami operator of
\Cref{def:geo:laplace-beltrami}. It is the manifold counterpart of the spectral
graph convolutions of \Cref{ch:graphs}, and it makes the analogy between the
continuous and discrete worlds exact.

\paragraph{The manifold Fourier transform.}
On a compact manifold $\Delta_\M$ has a discrete spectrum
$0 = \lambda_0 \le \lambda_1 \le \cdots$ with eigenfunctions $\{\varphi_k\}$
forming an orthonormal basis of $L^2(\M)$ (\Cref{rem:geo:lb-role}). These play
the role of complex exponentials: in $\R^n$ the Fourier modes $e^{i\inner{\xi}{x}}$
are exactly the eigenfunctions of the Euclidean Laplacian, and on a general
manifold the $\varphi_k$ generalize them. One defines the \emph{manifold Fourier
transform} and its inverse by
\[
  \hat f(\lambda_k) = \inner{f}{\varphi_k} = \int_\M f\, \varphi_k\, dV_g,
  \qquad
  f = \sum_{k} \hat f(\lambda_k)\, \varphi_k .
\]
Low eigenvalues correspond to smooth, slowly varying modes and high eigenvalues
to oscillatory ones, exactly as low and high frequencies in classical Fourier
analysis.

\paragraph{Spectral convolution and isometry invariance.}
Since there is no translation on a general manifold, convolution is
\emph{defined} through the convolution theorem: a spectral filter acts by
modulating each Fourier coefficient,
\begin{equation}
\label{eq:geo:spectral-conv}
  (g_\eta \star f) = \sum_k \eta(\lambda_k)\, \hat f(\lambda_k)\, \varphi_k,
\end{equation}
where $\eta$ is a learnable transfer function of the eigenvalue. Because
$\Delta_\M$ commutes with isometries, so does any filter of the form
\eqref{eq:geo:spectral-conv}: the architecture is automatically invariant under
the symmetries of the manifold, realizing the equivariance principle of
\Cref{ch:priors} with no extra constraints. Choosing $\eta$ to be a polynomial
in $\lambda$ makes the filter \emph{local} (a degree-$K$ polynomial reaches only
the $K$-hop geodesic neighbourhood) and avoids ever computing the full
eigendecomposition --- precisely the localization argument that, in the discrete
setting, produced Chebyshev and graph convolutional filters in \Cref{ch:graphs}.

\paragraph{The continuum--discretization bridge.}
The spectral construction is what makes the manifold and graph pictures two faces
of one theory. \Cref{tab:geo:analogy} lines up the correspondence: the graph
Laplacian discretizes $\Delta_\M$, its eigenvectors discretize the
eigenfunctions $\varphi_k$, and the graph and manifold Fourier transforms agree
in the sampling limit. A spectral filter learned on a sufficiently dense mesh
therefore approximates a continuous, mesh-independent operator --- the property
that lets such models transfer across discretizations of the same shape.

\begin{table}[h]
\centering
\caption{Correspondence between manifold and graph spectral analysis.}
\label{tab:geo:analogy}
\begin{tabular}{@{}p{5.4cm}p{6.0cm}@{}}
\toprule
\textbf{Manifold $(\M, g)$} & \textbf{Graph $G = (V,E)$} \\
\midrule
Laplace--Beltrami $\Delta_\M$ & Graph Laplacian $L = D - W$ \\
\addlinespace
Eigenfunctions $\varphi_k$ & Laplacian eigenvectors $u_k$ \\
\addlinespace
Spectrum $0 = \lambda_0 \le \lambda_1 \le \cdots$ & Eigenvalues $0 = \lambda_0 \le \cdots \le \lambda_{n-1}$ \\
\addlinespace
$\hat f(\lambda_k) = \inner{f}{\varphi_k}$ & $\hat f = U^\top f$ \\
\addlinespace
Spectral filter $\eta(\lambda_k)$ & Spectral filter $\eta(\Lambda)$ \\
\addlinespace
Volume element $dV_g$ & Vertex weights / Voronoi areas \\
\bottomrule
\end{tabular}
\end{table}

\paragraph{Outlook.}
Geodesics complete the geometric picture begun with grids, groups, and graphs:
they endow a domain with intrinsic distance, curvature, and a principled notion
of information flow. Two threads carry forward. The gauge ambiguity of normal
coordinates (\Cref{rem:geo:gauge}) --- the absence of a canonical local frame on
a curved domain --- is the precise problem solved by the gauge-equivariant
networks of \Cref{ch:gauges}. And the shared spectral language of
\Cref{tab:geo:analogy}, together with the pseudo-coordinate aggregation of
MoNet, places convolutions on grids, graphs, and manifolds inside the single
message-passing framework that \Cref{ch:applications,ch:conclusion} return to.
